% !TeX root = ../papers/nc-p3.tex
\section{三十節点中心のBrauer群と窓の八重被覆}
\label{sec:thirty-node-dihedral-cover}
\begin{theorem}\label{thm:thirty-node-dihedral-cover}
十五次元成分の一般中心 $X$ の不分岐Brauer群は巡回群 $\mathbb Z/4$ である。
中心の三十点のうちどの二十点に薄い単純があるかは $X$ だけで決まる。
この成分から中心四次式のモジュライへの写像は、一般点上で二面体群（位数8）のGalois被覆である。
同じ一般中心を持つ八つの添字保持同型類は、一つの窓とその反対窓の各四移動で尽くされる。
従ってこの開集合で、二つの窓の中心が射影同型であることは、添字の移動と逆転を許した同型と同値である。
一般中心に非自明な射影自己同型はない。
\end{theorem}
\begin{proof}[解消上のねじれ圏]
$Y$ を前節の三十点の吹上げとし、$\pi:Y\to X$ とする。まず複素数上で証明する。
前節の類を $Y$ に延ばし、その分裂のgerbeと、$U$ 上の折り畳み代数層
$\mathcal B|_U=\mathcal End(T_U)$ の重み1の階数16束を取る。
この束を $Y$ 上のねじれ束 $T$ に次のように延長できる。
各節点の完備局所模型を解析エタール局所で取れることを明記する。
節点は正則な二次主要部から解析的に通常二重点と同型である。標準模型と代数層との
完備同型は、有限個の加群生成元の像・逆像、関係と積と両合成の条件を満たす。
これらは有限表示の方程式にできる（加群の関係の係数も変数にする）。
Artin近似により解析局所環上で同型と逆を取れ、近傍を縮めても同じである。

標準節点の吹上げは $Q=\mathbb P^1\times\mathbb P^1$ 上の
$\mathcal O(-1,-1)$ の全空間の零切断近傍で、$\mathcal L$ を一方の罫線の
$\mathcal O(1,0)$ の引き戻しとする。$Q$ の小近傍ではgerbeを自明化できる。
実際この近傍は $Q$ に収縮し、$H^3(Q,\mathbb Z)=0$ である。
またSteinな節点近傍の逆像で $H^2(\mathcal O)=0$ である。
これは $Q$ の厚みの商 $\mathcal O_Q(n,n)$ の高次コホモロジーの消滅と
固有形式関数定理による。指数列を適用すればよい。
そこで節点型に応じて $\mathcal O^a\oplus\mathcal L^b$ を取る。
例外を除いた所では $T_U$ とは直線捩りだけ異なる（同じ自己準同型環の分裂束である）。
その局所的な直線捩りに任意の解析的関数を仮定する必要はない。
実際穴あき節点の小近傍はリンクに収縮し、リンクの $H^2(\mathbb Z)=\mathbb Z$ は
$Q$ 上の二つの次数を $(1,1)$ で割ったものである（円周束のGysin列）。
穴あき近傍では $H^1(\mathcal O)=0$ が深さ3とStein性から従うので、
指数列により全ての直線束はこの罫線束の冪である。従ってねじりを調整して
$T_U$ と同型にでき、互いに交わらない三十近傍で貼り合わせられる。
射影的 $Y$ 上のこの解析的ねじれ束は代数的である。例えば $\alpha$ を表す
Azumaya代数の一つで通常の加群層に移してGAGAを使い、作用写像も一緒に代数化する。

$R\pi_*\mathcal End(T)=\mathcal B$ である。局所標準模型では二つの罫線の一方の
小解消へ吹き下ろすと $\mathcal O\oplus\mathcal L$ はそこにある束の引き戻しである。
小解消では余次元2の例外曲線を越えて束準同型が全て延びる。従って零次直像は
穴あき節点上の自己準同型の反射的延長に等しい。
高次直像の消滅は例外二次曲面とその厚み上で
$\mathcal O_Q(n+j,n)$、$n\ge0$, $j=0,\pm1$ を調べると従う。
これから $T$ を核とする関手
$\operatorname{Perf}(X,\mathcal B)\to D^b(Y,\alpha)$ は充満忠実である。
右随伴は $F\mapsto R\pi_*\mathcal Hom(T,F)$（導来Hom）で、単位が同型であることは
局所射影上のこの直像計算と射影分解による。右随伴の像は局所大域次元の有限性で完全である。
左辺には命題\ref{prop:extra-central-singular-characters}の証明で得た長さ4の充満例外列がある。

各例外 $j:Q\hookrightarrow Y$ で、選んだ自明化のもと
$j_*\mathcal O_Q(-1,-1),j_*\mathcal O_Q(0,-1)$ を加える
（共通直線捩りを $T$ の局所規格化に合わせる）。
この二対象は先の像からの全Extが零である。$T|_Q$ の二つの成分との
Hom束では第二次数が $-1$ になるからである。
またこの順で一つの例外対である。除数の自己Extには $Q$ 上でのExtと
法束 $\mathcal O_Q(-1,-1)$ を掛けた一段低いExtだけが現れ、後者は同じ理由で零である。
異なる例外の間には射がない。従ってこれらの対を先に、先の像を後に置けば例外列となる。

生成性も局所的に確認する。この全ての対象からのExtが零な $F$ は、随伴により
$R\pi_*\mathcal Hom(T,F)=0$ なので例外上に台を持つ。
各完備局所節点への平坦基底変換後、$\mathcal O,\mathcal L$ からの全Extが零である。
例外対からの消滅も、台上のExtの完備化により保たれる。
例外除数の列
\[
 0\to\mathcal O(t+1,0)\to\mathcal O(t,-1)
        \to j_*\mathcal O_Q(t,-1)\to0\qquad(t=-1,0)
\]
により、$\mathcal O(-1,-1),\mathcal O(0,-1),\mathcal O,\mathcal O(1,0)$ の
四直線束からの全Extが零である。これらは $Q$ 上の生成系の引き戻しである。
実際第二因子の $\mathcal O(-1),\mathcal O$ の各層で第一因子の連続する二直線束を取っている。
完備節点の吹上げから $Q$ への射はアフィンである。標準アフィン節点の
吹上げは上の直線束の全空間であり、完備局所環へのアフィン基底変換を合成すればよい。
従ってこの四対象からの消滅は $F$ の直像、従って $F$ 自身の消滅を意味する。
これで $D^b(Y,\alpha)$ は長さ $4+2\cdot30=64$ の充満例外列を持つ。
\end{proof}

\begin{proof}[位相的Brauer類と節点の区別]
この充満性から、位相的なねじれ複素K群 $K^1(Y,\alpha)=0$ が従う。
理由を記す。対角の恒等核から各例外対象への評価・射影核を順に外すと、
最後の核は全対象上で零、特に各点の層上で零なので零である。
評価核は例外対象の双対と例外対象の外積であり、それらの合成も外積の有限複体になる。
従って対角の類は逆ねじれと順ねじれの外積の整数線型結合である。
位相的K群に移して各外積対応を作用させると、奇数次の入力との積の指数は
点の奇数K群に入って零となる。恒等対応も零となり主張が従う。

ねじれ複素K理論のAtiyah--Hirzebruchスペクトル系列で $H^3(Y,\mathbb Z)$ に
入る微分は $d_3:H^0\to H^3$ だけで、$1$ の像は符号を除き $\alpha_{\rm top}$ である。
ここで $d_3$ は整数Steenrod作用 $\operatorname{Sq}^3_{\mathbb Z}$ とねじれ類との
カップ積の和である。$H^3$ から $H^6$ へのその像は捩れであり、向き付けられた
コンパクト六次元多様体の最高コホモロジーは自由なので零である。他の出入りは次元上ない。
従って $H^3(Y,\mathbb Z)/\langle\alpha_{\rm top}\rangle$ は $K^1$ の随伴次数部分で、零である。
$H^2(Y,\mathcal O_Y)=0$ と指数列・有限係数の比較から
$\operatorname{Br}(Y)$ は $H^3(Y,\mathbb Z)$ の捩れ部分と同型である。
前節の周期計算によりこの群は $\mathbb Z/4$ である。$U$ 上のどのBrauer類も
穴あき完備節点上では零であり、系\ref{cor:quartic-center-unramified-brauer}と同じ
剰余の議論で $Y$ に延びる。従って $\operatorname{Br}(U)$ もこの群に一致する。
標数零の代数閉体の間の拡大では、滑らかな射影解消の有限係数エタールコホモロジーと
N\'eron--Severi群が変わらないので同じ結論である。

$\alpha$ または $-\alpha$ を選び、$H^2(U,\mu_4)$ に持ち上げる。
穴あき完備節点への制限は局所Picard群を4で割った元になる。
持ち上げの違いは大域の直線束であり、その局所類は前節で零と示した。
この元は符号を除き局所分裂後の階数4束の行列式で測れる。
実際重み4の直線束 $\det E_i$ を自明化する重み1の四乗根の障害がこのKummer持ち上げであり、
局所で重み1の自明化を変えても四倍の差しかない。二十点では奇数で、
残り十点では二種類とも頂点次元2のため $2\bmod4$ である。
生成元を取り替えても符号だけなので、これで順序なしの中心自身が二集合を区別する。
\end{proof}

\begin{proof}[八枚の区別]
向きを逆転すれば自己準同型環が反対になり、中心は同じでBrauer類は負になる。
開始位置の変更は類を保つ。この二操作は十五次元成分を保つ。
実際、定理\ref{thm:split-surface-component}の証明に使った位数4の量子窓では、
五次元以上の四次中心を持つ局所軌跡がまさにこの成分であり、この条件は両操作で不変である。
量子窓には周期1があり、その逆交換係数の三角積の組も元の組と文字の置換で移り合う
（四三角を向き付けた値は全て等しく、奇置換で反転する）。
従って両操作はその局所成分、よって既約成分全体を保つ。

$X$ と同定した中心の上では、指数4の類は上の二生成元だけである。
定理\ref{thm:fixed-brauer-center-windows}により各向きで四枚が全てである。
さらに中心同定を忘れてもこの八枚は異なる。定理\ref{thm:full-fifteen-component}の証明で、
一般窓の四位置の二十点配置は互いに射影非同値である。
非自明な巡回移動も、四頂点の逆転のどの移動も少なくとも一対の位置を
入れ替えるため、対応する添字保持同型は存在しない。
従って窓の一般点に自由に作用する位数8の二面体群の軌道が中心写像の一般繊維そのものである。
写像は標数零で生成的有限分離なので、中心成分の関数体はこの群の不変体であり、
開集合を縮めればエタールなGalois被覆となる。
最後に中心自己同型で一枚の中心の印を動かすと、窓自体の添字保持同型類は変わらない。
八枚の区別により印付きでも同じで、従って窓自己同型がその中心自己同型を誘導する。
一般窓の射影安定群は自明で、残る文字倍率は中心を射影的に固定する。よって中心自己同型も自明である。
\end{proof}
